\input{preamble}

\NewDocumentCommand{\rrImplement}{m m}{%
	\ensuremath{\pm%
	\IfNoValueTF{#2}{#1}{\tfrac{#1}{#2}}%
	}%
}
\NewDocumentCommand{\rr}{> {\SplitArgument{1}{/}} m}{%
	\rrImplement #1%
}

\newcommand{\drs}[3]{%
%Positive: #1\\%
%$P(-x)=#2$\\%
%Negative: #3
Positive: #1; Negative: #3
}

\NewDocumentCommand{\SynDiv}{m >{\SplitArgument{4}{;}} m >{\SplitArgument{4}{;}} m >{\SplitArgument{4}{;}} m o}{%
\IfNoValueF{#5}{#5:\\}%
\begin{tikzpicture}[x=3.5em,y=-3.5ex,baseline=-3.5ex]
\useasboundingbox(0,0) rectangle (6,3);
\SynDivRow{1}#2
\SynDivRow{2}#3
\SynDivRow{3}#4
\SynDivNum{0}{1}{#1}
\begin{scope}[yshift=-1ex]
\draw (0,1) -- (0.75,1) -- (0.75,0);
\draw (0.75,2) -- ++(5,0);
\draw (4.75,2) -- ++(0,1);
\end{scope}
\end{tikzpicture}%
}
%\NewDocumentCommand{\SynDivRowHelper}{m >{\SplitArgument{4}{;}} m}{\SynDivRow{#1}#2}
\NewDocumentCommand{\SynDivRow}{m m m m m m}{%
\IfNoValueTF{#6}{
\foreach \x/\lbl in {2/#2,3/#3,4/#4,5/#5} {\SynDivNum{\x}{#1}{\lbl}}
}{
\foreach \x/\lbl in {1/#2,2/#3,3/#4,4/#5,5/#6} {\SynDivNum{\x}{#1}{\lbl}}
}
}
\newcommand{\SynDivNum}[3]{%
\draw (#1,#2) node[anchor=base,text width=3.5em,align=right]{$#3$};
}

\title{Section 5.3 --- Rational Roots Theorem --- Answer Key}

\begin{document}
\maketitle

\section*{Multiplying Quickly}
Find the leading term of each product below.
\begin{smenum}
\mitemxxx
{$2x^4$}
{$8x^6$}
{$12x^9$}
\end{smenum}
Find the constant term of each product below.
\begin{smenum}
\mitemxxx
{$28$}
{$28$}
{$96$}
\end{smenum}

\section*{Possible Rational Roots}
List all possible rational roots of each polynomial below.
\begin{smenum}
\mitemxxx
{\rr{1}, \rr{2}, \rr{3}, \rr{6}}
{\rr{1}, \rr{5}, \rr{1/3}, \rr{5/3}}
{\rr{1}, \rr{2}, \rr{1/2}, \rr{1/3}, \rr{2/3}, \rr{1/6}}
\mitemx
{\rr{1},   \rr{2},   \rr{3},   \rr{5}, \rr{6}, \rr{10},   \rr{15}, \rr{30},
 \rr{1/2},           \rr{3/2}, \rr{5/2},                  \rr{15/2},
 \rr{1/3}, \rr{2/3},           \rr{5/3},       \rr{10/3},
 \rr{1/4},           \rr{3/4}, \rr{5/4},                  \rr{15/4},
 \rr{1/6},                     \rr{5/6},
 \rr{1/12},                    \rr{5/12}}
\end{smenum}

\section*{Descartes's Rule of Signs}
Use Descartes's Rule of Signs to describe the possible numbers of positive and negative roots of each polynomial.
\begin{smenum}
\mitemxx
{\drs{2 or 0}{x^4+x^3-7x^2-x+6}{2 or 0}}
{\drs{1}{-x^5+2x-4}{2 or 0}}
\mitemxx
{\drs{0}{-x^3-x+1}{1}}
{\drs{5, 3, or 1}{-x^5-10x^4-35x^3-50x^2-24x-7}{0}}
\end{smenum}

\section*{Trial and Error}
For the polynomial $12x^4-23x^3-58x^2+37x+60$, test each of the following possible rational roots, and state which (if any) of the following are true:
\begin{clist}
\item It is a root
\item It gives an upper bound on the roots
\item It gives a lower bound on the roots
\item None of these
\end{clist}
\begin{smenum}
\mitemxx
{\SynDiv{1}{12;-23;-58;37;60}{12;-11;-69;-32}{12;-11;-69;-32;28}[None of these]}
{\SynDiv{-2}{12;-23;-58;37;60}{-24;94;-72;70}{12;-47;36;-35;130}[Lower bound]}
\mitemxx
{\SynDiv{3}{12;-23;-58;37;60}{36;39;-57;-60}{12;13;-19;-20;0}[Root]}
{\SynDiv{5}{12;-23;-58;37;60}{60;185;635;3360}{12;37;127;672;3420}[Upper bound]}
\mitemx
{\SynDiv{-\frac{4}{3}}{12;-23;-58;37;60}{-16;52;8;-60}{12;-39;-6;45;0}[Root]}
\end{smenum}

\section*{Factor Fully}
Each polynomial below can be factored as a product of linear factors with integer coefficients. Factor them. Begin by using Descartes's Rule of Signs to analyze the signs of the roots.
\begin{smenum}
\mitemxx
{$(x+1)(x+3)(x-2)$}
{$(x+2)(x+5)(2x-3)$}
\mitemxx
{$(x+1)(2x+1)(2x+5)(3x-2)$}
{$(x+1)(x+2)(x+3)(x-4)(x-5)$}
\end{smenum}

\end{document}